% Bagian: first-order-logic
% Bab: model-theory
% Subbagian: overspill

\documentclass[../../../include/open-logic-section]{subfiles}

\begin{document}

\olfileid[id]{mod}{bas}{ove}
\olsection{Peluapan}

\begin{thm}
\ollabel{overspill} Jika suatu himpunan kalimat $\Gamma$ mempunyai model
hingga dengan ukuran sebesar apa pun, maka himpunan itu mempunyai model tak
hingga.
\end{thm}

\begin{proof}
Ekspansikan bahasa dari $\Gamma$ dengan menambahkan terhitung banyak simbol
konstanta baru $c_0$, $c_1$, \dots, lalu tinjau himpunan $\Gamma \cup
\{c_i \neq c_j : i \neq j\}$. Mengatakan bahwa $\Gamma$ mempunyai model
hingga dengan ukuran sebesar apa pun berarti bahwa bagi setiap $m >0$
terdapat $n\ge m$ sedemikian sehingga $\Gamma$ mempunyai model berkardinalitas
$n$. Hal ini menyiratkan bahwa $\Gamma \cup \{c_i \neq c_j : i \neq j\}$
dapat dipenuhi secara berhingga. Menurut kekompakan, $\Gamma \cup \{c_i
\neq c_j : i \neq j\}$ mempunyai suatu model $\Struct M$ yang domainnya
harus tak hingga, karena model itu memenuhi semua pertidaksamaan $c_i \neq
c_j$.
\end{proof}

\begin{prop}
\ollabel{inf-not-fo}
Tidak ada kalimat $!A$ dari bahasa orde pertama mana pun yang benar dalam
!!a{structure}~$\Struct M$ jika dan hanya jika domain $\Domain{M}$ dari
!!{structure} tersebut tak hingga.
\end{prop}

\begin{proof}
Jika kalimat~$!A$ semacam itu ada, negasinya $\lnot !A$ akan benar tepat
dalam semua !!{structure}s hingga. Karena itu, negasi tersebut akan mempunyai
model hingga dengan ukuran sebesar apa pun, tetapi tidak mempunyai model tak
hingga, bertentangan dengan \olref{overspill}.
\end{proof}

\end{document}
